% This file contains the abstract part of your thesis - in English and
% in Hebrew (within \abstractEnglish and \abstractHebrew respectively).
%
% Notes:
% - This file uses the UTF-8 character set encoding for the Hebrew
%   text not to get garbled. Keep it that way.
% - Assuming your thesis is mainly in English, Graduate School
%   regulations mandate the following lengths for the abstracts:
%
%      Language    Min. Length   Max. Length
%     ---------------------------------------
%      English       200 words     500 words
%      Hebrew        500 words   2,000 words
%
%   so that the Hebrew abstract typically has some content from
%   the English introduction and an overview of the results, not
%   present in the English; it is not just a translation.

\abstractEnglish{

Modern CPUs rely on undocumented performance mechanisms, such as branch prediction, prefetching, and store-load reordering, that introduce hard-to-find vulnerabilities. Model-based relational testing (MRT) tools such as Revizor have made black-box fuzzing for these leaks practical on x86-64. Arm, which powers nearly all mobile devices and a growing share of HPC and datacenters, has received far less scrutiny, despite being vulnerable to Spectre-class and Arm-specific attacks such as PACMAN and TikTag.

Many of these leaks share a structure: an instruction imposes preconditions on its operands, and the handling of a violated precondition leaves an observable microarchitectural trace. We port Revizor to AArch64, add support for Pointer Authentication (PAuth) and the Memory Tagging Extension (MTE), enable headless fuzzing of remote cores, and introduce a method for fuzzing instruction preconditions. We also show that Revizor's validation procedure discards leaks whose source lies outside the L1D cache, and add a source-agnostic detection method that, while observing only the L1D, provably extends the set of detectable leaks to those carried by any microarchitectural structure whose effect can be made to reach the L1D.

We discover two new speculative leaks affecting both server-class and mobile Arm CPUs: one in speculative canonicality and alignment checks, and the other in the sign extension performed by Pointer Authentication instructions. We also identify a known hardware-prefetcher leak on x86-64 that Revizor could not detect previously. Both Arm leaks cross privilege (EL0-attacker to EL2-victim) and virtualization (attacker-VM to victim-host) boundaries. We turn the canonicality leak into a covert channel, recover a victim's stack canary, and, through a ROP chain in OpenSSL, extract a complete RSA-2048 key.

} % end of English abstract


\abstractHebrew{

% Note that certain commands don't work that well in Hebrew "mode".
% If this happens to you, try wrapping the command within a
% \textenglish{ } - that may (or may not) help.

מעבדים מודרניים מתבססים על מנגנונים רבים, שרובם אינם מתועדים, כדי לשפר את ביצועיהם. המעבד מהיר בכמה סדרי גודל מרכיבים אחרים כדוגמת הזיכרון, ולכן לעיתים קרובות הוא נדרש לבחור בין שתי אפשרויות. האפשרות הראשונה היא לעצור בכל פעם שהוא זקוק למידע שאין לו גישה מהירה אליו, ולהמתין לתשובה מהימנה. האפשרות השנייה היא לבצע ניחוש מושכל, ופעולה זו נקראת ספקולציה.

מעבדים מודרניים עושים שימוש בסוגים רבים של ספקולציה, ותרומתם לביצועים אינה מוטלת בספק. חיזוי הסתעפויות \textenglish{(branch prediction)}, למשל, מנסה לחזות את תוצאותיהן של הסתעפויות בקוד. סידור מחדש של פקודות \textenglish{store} ו-\textenglish{load} \textenglish{(store-load reordering)} מאפשר לפקודות \textenglish{load} להתבצע לפני פקודות \textenglish{store} שקודמות להן. לא ניתן לעשות זאת ללא ספקולציה תוך שמירה על נכונות ארכיטקטונית. מנגנון ה-\textenglish{hardware prefetching} מנסה לחזות לאילו כתובות בזיכרון התוכנית תיגש, עוד לפני שביקשה זאת במפורש.

אופטימיזציות מעין אלו חשובות מאוד, אך הן גם הובילו לפרצות אבטחה חמורות רבות. התקפה על מנגנונים אלו, בשילוב עם ערוץ צד כדוגמת זיכרון המטמון \textenglish{(cache)} או צריכת החשמל, עלולה להדליף מידע סודי. פרצות כדוגמת \textenglish{Spectre} ו-\textenglish{Meltdown}, שהתגלו לראשונה בארכיטקטורת \textenglish{x86}, הדגימו זאת. הן ניצלו מספר מנגנונים מעין אלו, השתמשו ב-\textenglish{cache} כערוץ צד, והצליחו להדליף מידע מעבר לגבולות אבטחה שהמעבד אמור לאכוף. דליפות חדשות ממשיכות להתגלות, על אף ניסיונות הגנה רבים.

מעבדי \textenglish{Arm} נמצאים ברוב המכריע של הטלפונים החכמים, והם הופכים לנפוצים יותר ויותר במחשבים ניידים ובחוות שרתים. למרות זאת, הם לא זכו לבחינה מעמיקה בהשוואה למקביליהם בארכיטקטורת \textenglish{x86-64}. מעבדי \textenglish{Arm} פגיעים להתקפות ספקולטיביות כמו \textenglish{Spectre}, וגם להתקפות ייעודיות כדוגמת \textenglish{TikTag} ו-\textenglish{PACMAN}.

דפוס חוזר עומד בבסיסן של דליפות ספקולטיביות רבות: פקודות מציבות תנאים על הקלט שלהן או על מצב המערכת כולה, ותנאים אלו נבדקים רק בזמן ריצה. למשל, תגיות \textenglish{MTE} נבדקות בזמן ביצוע פקודות גישה לזיכרון, ומצביעים עשויים לעבור אימות מול חתימת \textenglish{PAC}. המערכת מוודאת את נכונותם של ערכים אלו כדי להמשיך בהרצת התוכנית. לעומת זאת, הפרה של תנאים אלו עלולה להוביל לחריגות \textenglish{(faults)} או לאפקטים מיקרו-ארכיטקטוניים אחרים שתוקף יכול לצפות בהם, ובכך להדליף מידע סודי.

במחקר זה אנו מתאימים את \textenglish{Revizor} לארכיטקטורת \textenglish{AArch64}, מוסיפים יכולת לביצוע \textenglish{fuzzing} מרחוק, ומרחיבים את התמיכה בתכונות ייחודיות ל-\textenglish{Arm}, ובהן \textenglish{MTE} ו-\textenglish{PAuth}. בנוסף, אנו מזהים פער ביכולותיו הנוכחיות של \textenglish{Revizor}, המונע ממנו לזהות משפחה גדולה של התקפות ספקולטיביות. אנו מציעים שיטה חדשה ויעילה לאיתור דליפות, ומוכיחים כי היא מזהה את כל מה ש-\textenglish{Revizor} מסוגל לזהות כיום, ואף יותר. לבסוף, אנו מציגים מתודולוגיה לבדיקה שיטתית של התנאים שפקודות מציבות להרצה מוצלחת.

יחד, שתי התרומות האלו מאפשרות לנו לחשוף שתי דליפות חדשות במעבדי \textenglish{Arm}. הראשונה נובעת מבדיקות התקינות שמעבדי \textenglish{Arm} מבצעים על הכתובת בכל גישה לזיכרון (קריאה, כתיבה והרצה). השנייה נובעת מתיקון אבטחה שנועד להבטיח שפקודות האימות של \textenglish{PAuth} אינן מדליפות את נכונות החתימה. אנו מראים כי בפועל פקודות אלו מדליפות תמיד את ביט \textenglish{55} של המצביע. בנוסף, אנו מזהים התקפה ידועה על מעבדי \textenglish{x86-64} ש-\textenglish{Revizor} לא היה מסוגל למצוא קודם לכן.

אנו מדגימים את הדליפות על מעבדי \textenglish{Neoverse N3} ו-\textenglish{Cortex-A715}, המיועדים לשרתים ולמכשירים ניידים בהתאמה, ומראים שהן חוצות גבולות הרשאה וגבולות וירטואליזציה. לבסוף, אנו מציגים שתי התקפות המנצלות את הדליפה הראשונה. הראשונה מדליפה את ה-\textenglish{stack canary} הסודי של תוכנית הקורבן. השנייה משתמשת בהתקפת \textenglish{ROP} ייעודית שפיתחנו עבור \textenglish{OpenSSL}, ומדליפה מפתח \textenglish{RSA-2048} פרטי שלם.

} % end of Hebrew abstract
